% Formatvorlage für Abschlussarbeiten an der Universität Wien
%
% Diese Vorlage ist für die Einreichung von Masterarbeiten vorgesehen und basiert auf einer Vorlage von Martin Dockner. Sie entspricht den verbindlichen Vorgaben zur Titelblattgestaltung an der Universität Wien. 
%
% Fügen Sie in diese Vorlage mittels Copy & Paste die Daten ein, die Ihnen im Schritt 6 im Prozess zum Upload der Abschlussarbeit zur Plagiatsprüfung in u:space angezeigt werden. Verändern Sie die Daten bitte nicht. Ein fehlerhaftes Titelblatt kann das Abschlussprozedere verzögern.  
%
% Version 0.9 2023
%
% Diese Vorlage nutzt die KOMA-Script-Klassen. Für Details siehe hier:
% ftp://ftp.dante.de/pub/tex/macros/latex/contrib/koma-script/scrguide.pdf
%%%%%%%%%%%%%%%%%%%%%%%%%%%%%%%%%%%%%%%%%%%%%%%%%%%%%%%%%%%%%%%%%%%%%%%%%%%%%%

\documentclass[12pt,pagesize,twoside,DIV=12,headsepline=0.4pt]{scrbook}

% Verwenden des Fonts Source Sans Pro
\usepackage[default,semibold,light]{sourcesanspro}
\usepackage[T1]{fontenc}


% \usepackage[rm={lining,proportional},sf={lining,proportional},tt={lining,tabular,monowidth}]{cfr-lm}
% \global\let\bfseries\sbweight



% folgende drei Pakete sind nur in deutsch wichtig! Sie können für englische
% Arbeiten auskommentiert werden ( % davor setzen)
\usepackage[english,ngerman]{babel}		% deutsche Abteilungsregeln
%\usepackage{lmodern} 	        % notwendig damit im PDF nach Umlauten gesucht werden kann
\usepackage[utf8]{inputenc}     % notwendig für Darstellung von Umlauten und scharfem S

% weitere wichtige Pakete
\usepackage{graphicx} 			% für Abbildungen benötigt
%\usepackage{natbib,bibgerm} 	% werden für die meisten Literaturstile benötigt
\usepackage[natbibapa]{apacite} % für Literaturestile nach den APA Zitierregeln
\usepackage{booktabs}			% für hübschere Tabellen
\usepackage{lineno}				% für Zeilnennummerierung
\usepackage{setspace}			% für 1.5-fachen oder doppelten Zeilenabstand
%\usepackage[DIV=15]{textarea}   % für Satzspiegeleinstellung/Layout

\usepackage{scrlayer-scrpage}   % für Modifkationen an Kopf- und Fusszeilen
    \pagestyle{scrheadings}     % völlig frei definierte Kopf- und Fusszeile
        \automark[chapter]{chapter} % nötig für den folgenden Befehl
        \ohead[]{\headmark}     % fügt in alle Kopfzeilen den Kapitelnamen ein
        \cfoot[]{}              % die Mitte der Fusszeile bleibt leer!
        \ofoot[\pagemark]{\pagemark} % fügt Seitenzahlen in der Fusszeile ein

% folgende Zeilen aktivieren Serifenschrift im kompletten Dokument
\setkomafont{disposition}{\normalcolor\bfseries}		% für Überschriften
\setkomafont{descriptionlabel}{\normalcolor\bfseries}	% für Descriptions
\setkomafont{captionlabel}{\normalcolor\bfseries}		% für Captions


%%%%%%%%%%%%%%%%%%%%%%%%%%%%%%%%%%%%%%%%%%%%%%%%%%%%%%%%%%%%%%%%%%%%%%%%%%%%%%
% Definition der Persönlichen Daten %%%%%%%%%%%%%%%%%%%%%%%%%%%%%%%%%%%%%%%%%%
%%%%%%%%%%%%%%%%%%%%%%%%%%%%%%%%%%%%%%%%%%%%%%%%%%%%%%%%%%%%%%%%%%%%%%%%%%%%%%
% Persönliche Angaben zur Arbeit folgen hier

\renewcommand{\title}{Zusatz zum Titel Titel meiner Arbeit} % Titel der Arbeit
\newcommand{\akademischegrade}{}            % bereits erlangte akademische Grade!
\renewcommand{\author}{Mag. Max Mustermann, BA}      % Titel Vor Vorname Nachname, Titel Nach
\newcommand{\matrikelnummer}{00000000}      % Matrikelnummer
\newcommand{\studienkennzahl}{A 000 000}                % Studienkennzahl
\newcommand{\studienrichtung}{Masterstudium Beispiel}    % Studienrichtung
\newcommand{\betreuerin}{ao. Univ. Prof. Dr. Max Mustermann, MSc}        % BetreuerIn

% Wenn du das Datum ändern willst, entferne in der nächsten Zeile
% Das Prozentzeichen und paße das Jahr händisch an.
% \year=2013
%%%%%%%%%%%%%%%%%%%%%%%%%%%%%%%%%%%%%%%%%%%%%%%%%%%%%%%%%%%%%%%%%%%%%%%%%%%%%%

\begin{document}						% Hier fängt der Text an
%%%%%%%%%%%%%%%%%%%%%%%%%%%%%%%%%%%%%%%%%%%%%%%%%%%%%%%%%%%%%%%%%%%%%%%%%%%%%%
% Hier wird die Titelseite eingefügt. Sie befindet sich in der Datei
% "titlepage.tex" und entspricht den Vorgaben der Uni Wien.
% Fügen Sie in diese Vorlage mittels Copy & Paste die Daten ein, die Ihnen im Schritt 6 im Prozess zum Upload der Abschlussarbeit zur Plagiatsprüfung in u:space angezeigt werden. Verändern Sie die Daten bitte nicht.
\include{titlepage}
% Ende der Titelseite
%%%%%%%%%%%%%%%%%%%%%%%%%%%%%%%%%%%%%%%%%%%%%%%%%%%%%%%%%%%%%%%%%%%%%%%%%%%%%%
% Es folgen einige Hilfsmittel, die du einbinden kannst, aber nicht mußt!
% Zum verwenden, entferne einfach in der entsprechenden Zeile das Prozentzeichen

%\pagewiselinenumbers				% Seitenweise Zeilennummerierung einschalten
%\onehalfspacing					% aktiviert 1.5-fachen Zeilenabstand
%\doublespacing						% aktiviert doppelten Zeilenabstand

%%%%%%%%%%%%%%%%%%%%%%%%%%%%%%%%%%%%%%%%%%%%%%%%%%%%%%%%%%%%%%%%%%%%%%%%%%%%%%
\tableofcontents					% Inhaltsverzeichnis
%%%%%%%%%%%%%%%%%%%%%%%%%%%%%%%%%%%%%%%%%%%%%%%%%%%%%%%%%%%%%%%%%%%%%%%%%%%%%%
% Aus Datenschutzgründen wird der Text Ihrer Abschlussarbeit in pseudonymisierter Form zur Plagiatsprüfung weitergeleitet. Eine etwaige Danksagung muss sich innerhalb der ersten 6 Seiten der Arbeit befinden, da diese vor der Übermittlung an die Texterkennungssoftware entfernt werden.

\chapter*{Danksagung}
\label{ch:danksagung}
Ich bedanke mich besonders bei Martin Dockner für die Einführung in \LaTeX\ und die Dokumentenvorlage für diese Masterarbeit.

%%%%%%%%%%%%%%%%%%%%%%%%%%%%%%%%%%%%%%%%%%%%%%%%%%%%%%%%%%%%%%%%%%%%%%%%%%%%%%
% Die Zusammenfassung/Der Abstract scheint nicht im Inhaltsverzeichnis auf.
%%%%%%%%%%%%%%%%%%%%%%%%%%%%%%%%%%%%%%%%%%%%%%%%%%%%%%%%%%%%%%%%%%%%%%%%%%%%%%
\chapter*{\abstractname}

Hier gehört die Zusammenfassung in Deutsch hin.

\selectlanguage{english} 		% Ab hier gelten englische Abteilungsregeln
\vspace{1.3em} 					% um den Abstand zum Text richtig zu stellen
\section*{\huge \abstractname} 	% Trick! Abstract auf gleicher Seite.
\vspace{1.3em} 					% um den Abstand zum Text richtig zu stellen

Write your english abstract here.

\selectlanguage{ngerman} 		% Ab hier gelten wieder deutsche Abteilungsregeln
%%%%%%%%%%%%%%%%%%%%%%%%%%%%%%%%%%%%%%%%%%%%%%%%%%%%%%%%%%%%%%%%%%%%%%%%%%%%%%

%%%%%%%%%%%%%%%%%%%%%%%%%%%%%%%%%%%%%%%%%%%%%%%%%%%%%%%%%%%%%%%%%%%%%%%%%%%%%%
% Hier beginnt der Haupttext.
%%%%%%%%%%%%%%%%%%%%%%%%%%%%%%%%%%%%%%%%%%%%%%%%%%%%%%%%%%%%%%%%%%%%%%%%%%%%%%

\chapter{Fragestellung}
\label{ch:fragestellung} % für einen Textverweis auf das Kapitel: \ref{ch:fragestellung}

Die Fragestellung der Diplomarbeit. Warum? Wann? Wo? Wieso?
Zur Übung auch ein kurzes Blockzitat.

\begin{quote}
Dieser Text ist links und rechts Eingerückt, und im Blocksatz gesetzt. Um das zu beweisen, muss hier noch ein bißchen mehr Text stehen, als tatsächlich notwendig.
\end{quote}


%%%%%%%%%%%%%%%%%%%%%%%%%%%%%%%%%%%%%%%%%%%%%%%%%%%%%%%%%%%%%%%%%%%%%%%%%%%%%%
\chapter{Methodik}
\label{ch:methodik}	% für einen Textverweis auf das Kapitel: \ref{ch:methodik}

Hier folgt ein Beispiel für eine mathematische Formel in \LaTeX. Ursprünglich wurde \TeX\ (der Vorgänger von \LaTeX) von einem Mathematiker\footnote{als Beispiel einer Fußnote hier der Name des \TeX-Entwicklers: Donald Knuth} entwickelt, um solche Formeln schön setzen zu können.

\begin{displaymath}
a_{1} + a_{2} + \ldots + a_{\infty} = \sum_{1=1}^{\infty}a_{i}
\end{displaymath}


%%%%%%%%%%%%%%%%%%%%%%%%%%%%%%%%%%%%%%%%%%%%%%%%%%%%%%%%%%%%%%%%%%%%%%%%%%%%%%
\chapter{Ergebnisse}
\label{ch:ergebnisse}

Hier kommen die Ergebnisse der Arbeit hin. Im Folgenden übrigens ein Beispiel für das Einbinden einer Grafik

\begin{figure}[ht]
\centering
\includegraphics[width=13cm]{univie}
\caption[Kurzname Logo]{Das offizielle Logo der Universität Wien}
\label{fig:univie}		% für einen Textverweis mittels \ref{fig:univie}
\end{figure}


%%%%%%%%%%%%%%%%%%%%%%%%%%%%%%%%%%%%%%%%%%%%%%%%%%%%%%%%%%%%%%%%%%%%%%%%%%%%%%
\chapter{Diskussion}
\label{ch:diskussion}

Die Diskussion schließt die Arbeit ab. Im Folgenden gebe ich noch ein Beispiel für eine Tabelle, die unter Zuhilfenahme des Paketes \texttt{booktabs} hübschere Ergebnisse liefert, als die \LaTeX-üblichen Tabellen.

% Die folgende Tabelle basiert auf dem Paket {booktabs}. Sie sieht den
% Tabellen in echten Publikationen sehr ähnlich. Jedoch sind hier senkrechte
% Linien nicht ordentlich realisierbar (wie sie in Publikationen auch nie
% vorkommen). Um Tabellen mit senkrechten Linien zu erzeugen, empfiehlt sich
% weiterhin die im Kurs besprochene Lösung mit der Angabe \hline für horizontale
% Linien, statt den hier gezeigten \toprule, \midrule und \bottomrule Befehlen.

\begin{table}[ht]
\centering
\caption[Kurzname Tabelle]{Eine Tabelle zur Verdeutlichung von Tabellen in \LaTeX. Es handelt sich um eine hübschere Tabelle, als im Kurs besprochen, die dem Aussehen von Tabellen in echten Publikationen sehr nahe kommt.}
\begin{tabular}{clp{6cm}}
\addlinespace					% fügt einen kleinen Abstand hinzu
\toprule						% Top-Linie bei einer Tabelle
Buchstabe   &	Wort	&	Satz			\\
\midrule						% dünnere Linie für Tabellen
A			&	Affe	&	Der Affe sitzt auf dem Baum und traut sich nicht herunter	\\
B			&	Biene	&	Die Biene fliegt Kreise um ihren Stock, und ihr wird dennoch nicht schwindelig \\
K			&	Karibu	&	Das Karibu ist nicht wie du! \\
\bottomrule						% Schlußlinie bei einer Tabelle
\end{tabular}
\label{tab:Buchstaben}		    % für einen Textverweis mittels \ref{tab:Buchstaben}
\end{table}


%%%%%%%%%%%%%%%%%%%%%%%%%%%%%%%%%%%%%%%%%%%%%%%%%%%%%%%%%%%%%%%%%%%%%%%%%%%%%%
\chapter{Schlussfolgerung}
\label{ch:schlussfolgerung}

Die Conclusion sollte alles oben erarbeitete Zusammenfassen und zu einem Ergebnis bringen.

Zum Testen kommt hier noch ein Literaturzitat. Wie in \cite{Dockner2010} bereits anschaulich erklärt, bietet \LaTeX\ gerade für Master- oder PhD-Arbeiten enormes Potential.


%%%%%%%%%%%%%%%%%%%%%%%%%%%%%%%%%%%%%%%%%%%%%%%%%%%%%%%%%%%%%%%%%%%%%%%%%%%%%%

\listoffigures   	% Verzeichnis der Abbildungen
\listoftables     	% Verzeichnis der Tabellen

%\nocite{*} 		% ALLE Zitate aus der Literaturdatenbank werden übernommen!

% Der Stil der Literaturzitate. Folgende Möglichkeiten sind günstig.
%  	apacite		= Author(Datum) .. komplett englisch
%	gerapalike	= Author(Datum) .. komplett deutsch
%  	gerplain	= [1] 			.. komplett deutsch
%	apalike2	= Author(Datum) .. teilweise für deutsch geeignet
%  	apa 		= Author(Datum) .. komplett englisch
%	elsart-harv	= Author(Datum) .. komplett englisch
%	chicago		= Author(Datum) .. komplett englisch
%	plain		= [1]			.. komplett englisch
%	nature		= ²				.. komplett englisch

\bibliographystyle{apacite} 
\bibliography{literature}


%%%%%%%%%%%%%%%%%%%%%%%%%%%%%%%%%%%%%%%%%%%%%%%%%%%%%%%%%%%%%%%%%%%%%%%%%%%%%%

\end{document}						% Ende des Dokuments
